% =================
% Preface / Nota do Autor
% =================
\chapter*{Nota do Autor}
\addcontentsline{toc}{chapter}{Nota do Autor}

\lettrine[lines=3,findent=2pt,nindent=4pt]{\textcolor{sangue}{E}}{ste livro}
nasceu de fragmentos. Não foi escrito como uma narrativa linear, mas como um
arquivo --- memórias, laudos, prontuários, cartas e fotografias que, juntos,
contam a história de Joaquim Antônio Correia, o Paciente 1.260.

A estrutura em fragmentos reflete a própria natureza da memória: ela não vem
em ordem cronológica, mas em flashes. Em sons. Em cheiros. Em momentos que o
corpo guarda mesmo quando a mente tenta esquecer.

Por isso, o leitor é convidado a percorrer estes fragmentos na ordem que
preferir. Há um percurso tradicional --- do Sertão à Contaminação ---, mas
não é o único. Ao final de cada fragmento, as rotas apontam para onde se pode
seguir, com o número da página de cada destino; o mapa de rotas, mais
adiante, mostra todos os caminhos de uma só vez. Cada fragmento aponta para
outros, como as memórias apontam umas para as outras.

Alguns fragmentos são memórias de Joaquim. Outros são documentos --- laudos,
prontuários, cartas. Outros ainda são a investigação de Lúcia Mendes, a
psiquiatra forense que, quase meio século depois, tenta entender o que
aconteceu. E há os fragmentos de contaminação: esses pertencem a você, leitor.

Este livro é sobre o que acontece quando uma história não é contada. Quando o
trauma não é processado. Quando a dor de um homem se espalha como uma doença
--- atravessando décadas, atravessando corpos --- até encontrar um hospedeiro
novo.

Se você começar a sonhar com a vala depois de ler este livro, não se
surpreenda.

O diário do Paciente 1.260 está aberto.

\cleardoublepage
